% Part: methods
% Chapter: proofs
% Section: example-1

\documentclass[../../../include/open-logic-section]{subfiles}

\begin{document}

\olfileid[id]{mth}{prf}{ex1}

\olsection{Sebuah Contoh}

Contoh pertama kita adalah fakta sederhana berikut tentang gabungan
dan irisan himpunan. Contoh ini akan menggambarkan penguraian definisi,
pembuktian konjungsi dan klaim universal, serta pembuktian dengan kasus.

\begin{prop}
Untuk sebarang himpunan $A$, $B$, dan $C$, $A \cup (B \cap C) = (A
\cup B) \cap (A \cup C)$.
\end{prop}

Mari kita buktikan!{}

\begin{proof}
Kita ingin menunjukkan bahwa untuk sebarang himpunan $A$, $B$, dan
$C$, $A \cup (B \cap C) = (A \cup B) \cap (A \cup C)$.
\begin{quote}
Mula-mula kita menguraikan definisi ``$=$'' dalam pernyataan
proposisi. Ingat bahwa membuktikan dua himpunan identik berarti
menunjukkan bahwa kedua himpunan itu memiliki !!{element} yang sama.
Artinya, semua !!{element} dari $A \cup (B \cap C)$ juga merupakan
!!{element} dari $(A \cup B) \cap (A \cup C)$, dan sebaliknya. Bagian
``dan sebaliknya'' berarti bahwa setiap !!{element} dari $(A \cup B)
\cap (A \cup C)$ juga harus merupakan !!a{element} dari $A \cup (B
\cap C)$. Jadi, ketika menguraikan definisi, kita melihat bahwa kita
harus membuktikan suatu konjungsi. Mari kita catat hal ini:
\end{quote}
Menurut definisi, $A \cup (B \cap C) = (A \cup B) \cap (A \cup C)$
jika dan hanya jika setiap !!{element} dari $A \cup (B \cap C)$ juga
merupakan !!a{element} dari $(A \cup B) \cap (A \cup C)$, dan setiap
!!{element} dari $(A \cup B) \cap (A \cup C)$ merupakan !!a{element}
dari $A \cup (B \cap C)$.
\begin{quote}
Karena ini merupakan konjungsi, kita harus membuktikan setiap
konjungnya secara terpisah. Mari kita mulai dengan yang pertama: kita
akan membuktikan bahwa setiap !!{element} dari $A \cup (B \cap C)$
juga merupakan !!a{element} dari $(A \cup B) \cap (A \cup C)$.

Ini merupakan klaim universal, sehingga kita mengambil suatu
!!{element} sembarang dari $A \cup (B \cap C)$ dan menunjukkan bahwa
objek itu juga harus merupakan !!a{element} dari $(A \cup B) \cap (A
\cup C)$. Kita akan memilih sebuah variabel sebagai nama !!{element}
sembarang tersebut, misalnya~$z$. Bukti kita berlanjut:
\end{quote}
Pertama, kita buktikan bahwa setiap !!{element} dari $A \cup (B \cap
C)$ juga merupakan !!a{element} dari $(A \cup B) \cap (A \cup C)$.
Misalkan $z \in A \cup (B \cap C)$. Kita harus menunjukkan bahwa $z
\in (A \cup B) \cap (A \cup C)$.
\begin{quote}
Sekarang tiba waktunya untuk menguraikan definisi $\cup$ dan~$\cap$.
Misalnya, definisi $\cup$ adalah: $A \cup B = \Setabs{z}{z \in A
  \text{ atau } z \in B}$. Ketika kita menerapkan definisi itu pada
``$A \cup (B \cap C)$'', peran ``$B$'' dalam definisi kini dimainkan
oleh ``$B \cap C$'', sehingga $A \cup (B \cap C) =
\Setabs{z}{z \in A \text{ atau } z \in B \cap C}$. Jadi, asumsi kita
bahwa $z \in A \cup (B \cap C)$ sama dengan: $z \in
\Setabs{z}{z \in A \text{ atau } z \in B \cap C}$. Dan $z \in
\Setabs{z}{\dots z\dots}$ jika dan hanya jika \dots{} $z$ \dots,
yaitu, dalam kasus ini, $z \in A$ atau $z \in B \cap C$.
\end{quote}
Menurut definisi $\cup$, $z \in A$ atau $z \in B \cap C$.
\begin{quote}
Karena ini merupakan disjungsi, akan berguna untuk menerapkan
pembuktian dengan kasus. Kita mengambil kedua kasus dan menunjukkan
bahwa dalam masing-masing kasus, kesimpulan yang kita tuju (yaitu, ``$z
\in (A \cup B) \cap (A \cup C)$'') berlaku.
\end{quote}
Kasus 1: Andaikan $z \in A$.
\begin{quote}
Tidak banyak lagi yang dapat kita olah dari asumsi-asumsi kita. Jadi,
mari kita perhatikan apa yang dapat kita gunakan dalam kesimpulan.
Kita ingin menunjukkan bahwa $z \in (A \cup B) \cap (A \cup C)$.
Berdasarkan definisi $\cap$, untuk menunjukkan bahwa $z \in (A \cup
B) \cap (A \cup C)$, kita harus menunjukkan bahwa objek itu berada
baik dalam $(A \cup B)$ maupun dalam $(A \cup C)$. Namun, $z \in A \cup B$
jika dan hanya jika $z \in A$ atau $z \in B$, dan kita sudah memiliki
(sebagai asumsi kasus~1) bahwa $z \in A$. Dengan penalaran yang
sama---memakai $C$ sebagai pengganti $B$---kita memperoleh $z \in A \cup C$.
Argumen ini berjalan dalam arah terbalik, jadi mari kita catat
penalarannya dalam arah yang diperlukan oleh bukti.
\end{quote}
Karena $z \in A$, maka $z \in A$ atau $z \in B$, sehingga menurut
definisi~$\cup$, $z \in A \cup B$. Demikian pula, $z \in A \cup C$.
Namun, ini berarti bahwa $z \in (A \cup B) \cap (A \cup C)$, menurut
definisi~$\cap$.
\begin{quote}
Ini menuntaskan kasus pertama dalam pembuktian dengan kasus. Sekarang
kita ingin menurunkan kesimpulan dalam kasus kedua, yaitu ketika $z
\in B \cap C$.
\end{quote}
Kasus 2: Andaikan $z \in B \cap C$.
\begin{quote}
Sekali lagi, kita sedang berurusan dengan irisan dua himpunan. Mari
kita terapkan definisi~$\cap$:
\end{quote}
Karena $z \in B \cap C$, $z$ harus merupakan !!a{element} dari $B$
sekaligus dari $C$, menurut definisi~$\cap$.
\begin{quote}
Sekarang kita kembali melihat kesimpulan kita. Kita harus menunjukkan
bahwa $z$ berada baik dalam $(A \cup B)$ maupun dalam $(A \cup C)$.
Sekali lagi, penyelesaiannya langsung.
\end{quote}
Karena $z \in B$, maka $z \in (A \cup B)$. Karena $z \in C$, maka $z
\in (A \cup C)$ juga. Jadi, $z \in (A \cup B) \cap (A \cup C)$.
\begin{quote}
Di sini kita kembali menerapkan definisi $\cup$ dan $\cap$, tetapi
karena definisi-definisi tersebut sudah kita ingat kembali dan kita
sudah menunjukkan bahwa jika $z$ berada dalam salah satu dari dua
himpunan, maka ia berada dalam gabungan keduanya, kita tidak perlu
menyatakan langkah yang dilakukan secara begitu eksplisit.

Kita telah menuntaskan kasus kedua dalam pembuktian dengan kasus, jadi
sekarang kita dapat menyatakan kesimpulan pertama.
\end{quote}
Jadi, jika $z \in A \cup (B \cap C)$, maka $z \in (A \cup B) \cap (A
\cup C)$.
\begin{quote}
Sekarang kita tinggal menunjukkan arah sebaliknya, yaitu bahwa setiap
!!{element} dari $(A \cup B) \cap (A \cup C)$ merupakan !!a{element}
dari $A \cup (B \cap C)$. Seperti sebelumnya, kita membuktikan klaim
universal ini dengan mengasumsikan bahwa kita memiliki suatu
!!{element} sembarang dari himpunan pertama, lalu menunjukkan bahwa
objek itu harus berada dalam himpunan kedua. Mari kita nyatakan apa
yang akan kita lakukan.
\end{quote}
Sekarang, asumsikan bahwa $z \in (A \cup B) \cap (A \cup C)$. Kita
ingin menunjukkan bahwa $z \in A \cup (B \cap C)$.
\begin{quote}
Kini kita bekerja dari hipotesis bahwa $z \in (A \cup B) \cap (A \cup
C)$. Mudah-mudahan tidak terlalu membingungkan bahwa kita menggunakan
~$z$ yang sama seperti pada bagian pertama bukti. Ketika bagian itu
selesai, semua asumsi yang kita buat di sana tidak lagi berlaku,
sehingga sekarang kita dapat membuat asumsi baru mengenai $z$. Jika
hal ini membingungkan, gantilah $z$ dengan variabel lain dalam bagian
berikut.

Kita mengetahui bahwa $z$ berada baik dalam $A \cup B$ maupun dalam $A
\cup C$, menurut definisi~$\cap$. Menurut definisi $\cup$, kita dapat
menguraikan hal ini lebih lanjut menjadi: $z \in A$ atau $z \in B$,
dan juga $z \in A$ atau $z \in C$. Ini tampak seperti pembuktian
dengan kasus lagi---hanya saja kata ``dan'' membuatnya membingungkan.
Anda mungkin mengira bahwa ini berarti ada tiga kemungkinan: $z$
berada dalam $A$, $B$, atau $C$. Namun, itu keliru. Kita harus
berhati-hati, jadi mari kita telaah setiap disjungsi secara bergiliran.
\end{quote}
Menurut definisi $\cap$, $z \in A \cup B$ dan $z \in A \cup C$.
Menurut definisi $\cup$, $z \in A$ atau $z \in B$. Kita bedakan kasus.
\begin{quote}
Karena kita sedang berfokus pada disjungsi pertama, kita belum
memperoleh disjungsi kedua (dari penguraian $A \cup C$). Bahkan, kita
belum memerlukannya. Kasus pertama adalah $z \in A$, dan !!a{element}
dari suatu himpunan juga merupakan !!a{element} dari gabungan himpunan
itu dengan himpunan mana pun. Jadi, kasus~1 mudah:
\end{quote}
Kasus 1: Andaikan $z \in A$. Akibatnya, $z \in A \cup (B \cap C)$.
\begin{quote}
Sekarang untuk kasus kedua, $z \in B$. Di sini kita akan menguraikan
$\cup$ kedua dan melakukan pembuktian dengan kasus sekali lagi:
\end{quote}
Kasus 2: Andaikan $z \in B$. Karena $z \in A \cup C$, maka $z \in A$
atau $z \in C$. Kita bedakan kasus lebih lanjut:

Kasus 2a: $z \in A$. Maka, sekali lagi, $z \in A \cup (B \cap C)$.
\begin{quote}
Baiklah, ini agak aneh. Kita sebenarnya tidak memerlukan asumsi bahwa
~$z \in B$ untuk kasus ini, tetapi itu tidak menjadi masalah.
\end{quote}
Kasus 2b: $z \in C$. Maka $z \in B$ dan $z \in C$, sehingga $z \in B
\cap C$, dan akibatnya $z \in A \cup (B \cap C)$.
\begin{quote}
Ini mengakhiri kedua pembuktian dengan kasus, sehingga bagian kedua
telah selesai.
\end{quote}
Jadi, jika $z \in (A \cup B) \cap (A \cup C)$, maka $z \in A \cup (B
\cap C)$.
\end{proof}

\end{document}
